\PassOptionsToPackage{table,xcdraw}{xcolor}
\documentclass[11pt, a4paper]{article}

\usepackage[T1]{fontenc}
\usepackage[utf8]{inputenc}
\usepackage{tgpagella}
\usepackage[spanish, es-noshorthands]{babel}
\usepackage{amsmath, amssymb}
\usepackage{booktabs}
\usepackage{tabularx}
\usepackage[most]{tcolorbox}
\usepackage[american]{circuitikz}
\usepackage[margin=2cm]{geometry}
\usepackage{fancyhdr}

\pagestyle{fancy}
\fancyhf{}
\setlength{\headheight}{16pt}
\lhead{\textbf{UCB Sede Tarija} | IMT-221}
\chead{Suplemento Lab 3: Caracterización DC}
\rhead{Semana 8 - Sesión 2}
\lfoot{Circuitos Electrónicos II}
\rfoot{Página \thepage}
\renewcommand{\headrulewidth}{0.4pt}
\definecolor{ucbblue}{RGB}{0, 51, 102}

\begin{document}

\begin{center}
    \Large\textbf{\textcolor{ucbblue}{Suplemento Lab 3: Caracterización DC Ampliada del Amplificador Diferencial}}[cite: 13]
\end{center}

\begin{tcolorbox}[colback=yellow!10, colframe=ucbblue, title=\textbf{Aviso Importante}]
    \textbf{Este documento no reemplaza la guía anterior de Lab 3.[cite: 13]} Solo actualiza dos aspectos críticos[cite: 13]:
    \begin{enumerate}
        \item La alimentación nominal del LM741 pasa de $\pm15\,\text{V}$ a $\pm10\,\text{V}$[cite: 13].
        \item La prueba diferencial se amplía de un único punto a una pequeña caracterización DC completa[cite: 13].
    \end{enumerate}
\end{tcolorbox}

\section*{1. Propósito y Pregunta Central[cite: 13]}
La actividad debe permitir observar experimentalmente la siguiente cadena[cite: 13]:
\begin{equation*}
    \text{modelo ideal} \rightarrow \text{respuesta física} \rightarrow \text{error} \rightarrow \text{diagnóstico} \rightarrow \text{compensación} \rightarrow \text{validación} \text{[cite: 13]}
\end{equation*}
\textbf{Pregunta central:} ¿Qué tan bien reproduce un amplificador diferencial real la relación ideal $V_o=10(V_2-V_1)$, qué errores aparecen en DC y cuáles pueden reducirse mediante compensación?[cite: 13]

\section*{2. Actualización de Alimentación y Circuito[cite: 13]}
\begin{minipage}{0.45\textwidth}
    \raggedright % Evita el Underfull \hbox en párrafos estrechos
    \textbf{Nuevos parámetros de fuente[cite: 13]:}
    \begin{equation*}
        \boxed{V^+ = +10\,\text{V}, \quad V^- = -10\,\text{V}} \text{[cite: 13]}
    \end{equation*}
    Con referencia en $0\,\text{V}$[cite: 13].

    Antes de energizar[cite: 13]:
    \begin{itemize}
        \item Verificar fabricante y marcado del LM741[cite: 13].
        \item Verificar pinout[cite: 13].
        \item Confirmar que la alimentación es válida[cite: 13].
        \item Medir $+10\,\text{V}, 0\,\text{V}$ y $-10\,\text{V}$ con multímetro[cite: 13].
    \end{itemize}
\end{minipage}
\hfill
\begin{minipage}{0.5\textwidth}
    \begin{center}
        \begin{circuitikz}[scale=0.8, transform shape, thick]
            \draw (0,0) node[op amp, noinv input up] (opamp) {};
            \draw (opamp.+) to[R, l_={$R_3=10\,\text{k}\Omega$}, *-o] ++(-2.5,0) node[left]{$V_2$};
            \draw (opamp.-) to[R, l_={$R_1=10\,\text{k}\Omega$}, *-o] ++(-2.5,0) node[left]{$V_1$};
            \draw (opamp.+) -- ++(0,1.5) coordinate(tmp1) to[R, l={$R_4=100\,\text{k}\Omega$}] (tmp1 -| opamp.out) node[ground]{};
            \draw (opamp.-) -- ++(0,-1.5) coordinate(tmp2) to[R, l_={$R_2=100\,\text{k}\Omega$}] (tmp2 -| opamp.out) -- (opamp.out);
            \draw (opamp.out) to[short, *-o] ++(1,0) node[right]{$V_o$};
            \draw (opamp.up) -- ++(0,0.5) node[above]{$+10\,\text{V}$};
            \draw (opamp.down) -- ++(0,-0.5) node[below]{$-10\,\text{V}$};
        \end{circuitikz}
    \end{center}
\end{minipage}

\section*{3. Derivación Explícita del Modelo Ideal[cite: 13]}
Para comprender el origen de la predicción, realizaremos el análisis paso a paso con los valores nominales adoptados: $R_1=R_3=10\,\text{k}\Omega$ y $R_2=R_4=100\,\text{k}\Omega$[cite: 13].

Asumiendo un Op-Amp ideal en régimen lineal ($i_+ = i_- = 0$, $V_+ = V_-$):

\textbf{Paso 1: Tensión en la entrada no inversora ($V_+$).} Al ser $i_+ = 0$, $R_3$ y $R_4$ forman un divisor de tensión exacto:
\begin{equation*}
    V_+ = V_2 \frac{R_4}{R_3 + R_4} = V_2 \frac{100\,\text{k}\Omega}{10\,\text{k}\Omega + 100\,\text{k}\Omega} = V_2 \frac{100}{110} = V_2 \frac{10}{11}
\end{equation*}

\textbf{Paso 2: Condición de realimentación.} $V_- = V_+ = V_2 \frac{10}{11}$.

\textbf{Paso 3: Ley de Corrientes de Kirchhoff (KCL) en el nodo inversor ($V_-$).} Como $i_- = 0$:
\begin{equation*}
    \frac{V_1 - V_-}{10\,\text{k}\Omega} + \frac{V_o - V_-}{100\,\text{k}\Omega} = 0
\end{equation*}
Multiplicando toda la ecuación por $100\,\text{k}\Omega$ para eliminar denominadores:
\begin{align*}
    10(V_1 - V_-) + 1(V_o - V_-) &= 0 \\
    10V_1 - 10V_- + V_o - V_- &= 0 \\
    V_o &= 11V_- - 10V_1
\end{align*}

\textbf{Paso 4: Sustitución de $V_-$ y resultado final.} Sustituimos $V_- = V_2 \frac{10}{11}$:
\begin{equation*}
    V_o = 11\left(V_2 \frac{10}{11}\right) - 10V_1 = 10V_2 - 10V_1 = 10(V_2 - V_1)
\end{equation*}
Definiendo el voltaje diferencial como $V_d = V_2 - V_1$[cite: 13], llegamos a la relación ideal oficial:
\begin{equation*}
    \boxed{V_{o,\mathrm{ideal}} = 10V_d} \text{[cite: 13]}
\end{equation*}

\section*{4. Caracterización Diferencial DC[cite: 13]}
Mantenga $V_1 = 0\,\text{V}$ y varíe $V_2$[cite: 13]. \textbf{Importante:} Mida siempre los valores reales en los pines de entrada[cite: 13]. $V_d = V_{2,\mathrm{meas}} - V_{1,\mathrm{meas}}$[cite: 13]. $e = V_{o,\mathrm{meas}} - V_{o,\mathrm{ideal}}$[cite: 13].

\begin{table}[h]
    \centering
    \footnotesize % Evita el Overfull \hbox en la tabla
    \renewcommand{\arraystretch}{1.5}
    \begin{tabularx}{\textwidth}{|c| *{6}{>{\centering\arraybackslash}X|} }
    \hline
    \textbf{Pt.} & \textbf{$V_{1,\mathrm{meas}}$} & \textbf{$V_{2,\mathrm{meas}}$} & \textbf{$V_d$} & \textbf{$V_{o,\mathrm{ideal}}$} & \textbf{$V_{o,\mathrm{meas}}$} & \textbf{Error ($e$)} \\ \hline
    D0[cite: 13] & $0\,\text{V}$ (Nom)[cite: 13] & $0\,\text{mV}$ (Nom)[cite: 13] & & $0.00\,\text{V}$[cite: 13] & & \\ \hline
    D1[cite: 13] & $0\,\text{V}$ (Nom)[cite: 13] & $20\,\text{mV}$ (Nom)[cite: 13] & & $0.20\,\text{V}$[cite: 13] & & \\ \hline
    D2[cite: 13] & $0\,\text{V}$ (Nom)[cite: 13] & $50\,\text{mV}$ (Nom)[cite: 13] & & $0.50\,\text{V}$[cite: 13] & & \\ \hline
    D3[cite: 13] & $0\,\text{V}$ (Nom)[cite: 13] & $100\,\text{mV}$ (Nom)[cite: 13] & & $1.00\,\text{V}$[cite: 13] & & \\ \hline
    \end{tabularx}
\end{table}

\textbf{Interpretación de Datos[cite: 13]:}
\begin{enumerate}
    \setlength{\itemsep}{0pt}
    \item ¿La salida aumenta de forma lineal con $V_d$?[cite: 13]
    \item ¿La ganancia experimental (pendiente) parece cercana a 10?[cite: 13]
    \item ¿Existe salida distinta de cero cuando $V_d \approx 0$ (intercepto)?[cite: 13] ¿Qué mecanismo podría causarlo?[cite: 13]
\end{enumerate}
\textit{Guía de análisis:} Un offset constante apunta a un error DC referido a entrada. Una pendiente distinta de 10 apunta a un desajuste en las razones resistivas (mismatch)[cite: 13].

\section*{5. Prueba de Condición Común (Si el tiempo lo permite)[cite: 13]}
\begin{table}[h]
    \centering
    \footnotesize
    \renewcommand{\arraystretch}{1.5}
    \begin{tabularx}{\textwidth}{|c|>{\centering\arraybackslash}X|>{\centering\arraybackslash}X|>{\centering\arraybackslash}X|>{\centering\arraybackslash}X|}
    \hline
    \textbf{Pt.} & \textbf{$V_{1,\mathrm{meas}}$ (Nominal)} & \textbf{$V_{2,\mathrm{meas}}$ (Nominal)} & \textbf{$V_{o,\mathrm{meas}}$} & \textbf{Observación ($V_{o,\mathrm{ideal}} = 0\,\text{V}$)[cite: 13]} \\ \hline
    C0[cite: 13] & $0\,\text{V}$ (Nom)[cite: 13] & $0\,\text{V}$ (Nom)[cite: 13] & & \\ \hline
    C1[cite: 13] & $0.5\,\text{V}$ (Nom)[cite: 13] & $0.5\,\text{V}$ (Nom)[cite: 13] & & \\ \hline
    C2[cite: 13] & $1.0\,\text{V}$ (Nom)[cite: 13] & $1.0\,\text{V}$ (Nom)[cite: 13] & & \\ \hline
    \end{tabularx}
\end{table}

\section*{6. Compensación y Repetición[cite: 13]}
Verifique el circuito. Bajo $V_1=V_2=0$, registre $V_{o,\mathrm{before}}$[cite: 13]. Aplique el potenciómetro a los pines OFFSET NULL según datasheet[cite: 13]. Registre $V_{o,\mathrm{after}}$ y calcule $\Delta|e|=|e_{\mathrm{before}}|-|e_{\mathrm{after}}|$[cite: 13].

\textbf{Remedición:} Repita los puntos D0, D2 y D3.
\begin{itemize}
    \setlength{\itemsep}{0pt}
    \item ¿Qué cambió después de compensar? ¿Cambió el offset, la pendiente, o ambos?[cite: 13]
    \item ¿Qué error residual se puede explicar por mismatch en lugar de offset interno del Op-Amp?[cite: 13]
\end{itemize}

\section*{7. LTSpice - Actividad Suplementaria y Evidencia Mínima[cite: 13]}
Simular el circuito con $\pm10\,\text{V}$ bajo las condiciones D0-D3[cite: 13]. Si el tiempo no permite ambas, la regla es clara: $\boxed{\text{hardware físico} > \text{LTSpice}}$[cite: 13].

\textbf{Evidencia requerida:} Resistores medidos, $V_1, V_2$ medidos, tablas completas, error calculado, compensación aplicada y una interpretación clara de las discrepancias respecto al modelo ideal[cite: 13].

\end{document}
